\begin{remark*}
The ball theorems are now Mathlib-first. We work in the context
\[
  \texttt{variable \{X : Type u\} [MetricSpace X]}.
\]
An open ball is the set \texttt{Metric.ball x r}; a closed ball is the set
\texttt{Metric.closedBall x r}. Membership is rewritten with Mathlib's API:
\[
  \texttt{Metric.mem\_ball'}
  \qquad\text{and}\qquad
  \texttt{Metric.mem\_closedBall'}.
\]
The primed versions state membership using \texttt{dist x y}, matching the
usual hand notation \(d(x,y)\).
\end{remark*}

\begin{remark*}
The center containment theorem is:
\[
\begin{gathered}
  \texttt{theorem center\_mem\_ball}\\
  \texttt{\ \ \ \ (x : X) \{r : Real\} (radius\_positive : 0 < r) :}\\
  \texttt{\ \ \ \ x in Metric.ball x r := by}\\
  \texttt{\ \ exact Metric.mem\_ball\_self radius\_positive.}
\end{gathered}
\]
Mathlib already knows that the center lies in its own open ball when the radius
is positive. The proof is therefore a direct API call.
\end{remark*}

\begin{remark*}
For ball monotonicity, the proof exposes the hidden set-theoretic step. To
prove
\[
  \texttt{Metric.ball x r subset Metric.ball x s},
\]
Lean asks for an arbitrary point and its membership proof:
\[
\begin{gathered}
  \texttt{intro y point\_in\_smaller\_ball}\\
  \texttt{rw [Metric.mem\_ball'] at point\_in\_smaller\_ball ⊢}\\
  \texttt{exact point\_in\_smaller\_ball.trans\_le radius\_le.}
\end{gathered}
\]
After rewriting membership, the proof is the elementary implication
\[
  d(x,y) < r,\quad r \le s
  \quad\Longrightarrow\quad
  d(x,y) < s.
\]
\end{remark*}

\begin{remark*}
The theorem
\[
  \texttt{ball\_subset\_ball\_of\_mem}
\]
formalizes the standard geometric fact that open balls are open. If
\(y \in B(x,r)\), choose
\[
  \epsilon = r - d(x,y).
\]
The membership hypothesis says \(d(x,y) < r\), so \(\epsilon > 0\). For any
\(z \in B(y,\epsilon)\), the triangle inequality gives
\[
  d(x,z) \le d(x,y) + d(y,z)
    < d(x,y) + (r-d(x,y))
    = r.
\]
The Lean proof follows the same chain with \texttt{calc}, \texttt{dist\_triangle},
\texttt{linarith}, and \texttt{ring}.
\end{remark*}
